\chapter{Discussion and Limitations}
\label{ch:discussion}

This chapter interprets the experimental findings in the context of related work,
identifies the principal limitations of the approach, and maps out directions for
future improvement.

\section{What the Results Show About Sequence-Based Code Understanding}

The joint model achieves F1\,=\,0.931 on the synthetic held-out set --- a strong result
that places it in the range reported for fine-tuned CodeBERT and GraphCodeBERT baselines
on comparable binary classification benchmarks~\cite{fu2022linevul,ni2023diversevul}.
The semantic understanding tests nuance this headline figure.
The model passes all six variable-renaming cases and all six token-ablation cases,
demonstrating that its representations are not anchored to specific identifier strings
or isolated keywords.
These results are consistent with what one would expect from a model pre-trained on
hundreds of millions of code tokens: the attention mechanism has learned associations
between structural patterns --- user-controlled data flowing into an execution sink,
or an unchecked arithmetic expression in a size comparison --- that persist across
surface variations.

However, the model fails all four control-flow-dependent dead-code cases and one
data-flow-dependent minimal-fix case (CWE-78 C with a hardcoded system argument).
These failures reveal the ceiling of sequence-level encoding.
A transformer encoder operating on a token sequence can learn that certain patterns of
tokens co-occur in dangerous code, but it cannot determine whether a code path is
reachable, whether a variable's value at a given point is user-controlled, or whether
a pointer freed at one location is dereferenced at another.
Token attribution analysis (Section~\ref{sec:insight5}) makes this concrete:
for CWE-78, \texttt{sprintf} receives an attribution score of 0.671 on a structurally
safe function --- the highest score recorded across all test cases --- demonstrating that
the model has learned \texttt{sprintf} as an unconditional signal of danger rather than
attending to whether user-controlled data flows into the buffer.
By contrast, CWE-416 attribution correctly identifies \texttt{free} (score 0.506) and the
subsequently dereferenced pointer as the primary signals, consistent with a model that
has learned the use-after-free pattern structurally.
These are properties that require an explicit representation of program structure:
a control-flow graph for reachability, a data-dependence graph for taint analysis,
a points-to graph for alias resolution.
Pre-trained models that incorporate these structures --- such as
GraphCodeBERT~\cite{guo2021graphcodebert} and graph-augmented vulnerability detectors
such as Devign~\cite{zhou2019devign} and Reveal~\cite{chakraborty2022reveal} --- have
demonstrated gains on exactly the CWE classes where sequence models fail.
The results of this work motivate rather than undermine that line of research:
they show precisely which vulnerability classes would benefit from richer structural
representations.

\section{Cross-Language Transfer: Limits and Potential}

The ablation results and the per-CWE transfer analysis show that cross-language transfer
is highly uneven.
Where a vulnerability is expressed through the same surface API in both languages,
transfer is near-perfect.
Where the languages implement the same conceptual flaw through structurally different
mechanisms --- file permission management, TOCTOU detection --- transfer fails almost
entirely.

This finding has implications beyond this specific system.
It suggests that current pre-trained code models, despite being trained on multi-language
corpora, learn language-specific realisations of concepts rather than language-agnostic
representations.
A model that has seen Python's \texttt{os.chmod(path, 0o777)} does not automatically
learn to recognise C's \texttt{fchmod(fd, 0777)} as the same concept.
Achieving genuine cross-language transfer for structurally divergent CWEs would require
either explicit cross-language supervision --- paired examples of the same vulnerability
in both languages --- or a representation that operates at the semantic rather than
syntactic level.

The asymmetry between C-to-Python and Python-to-C transfer reflects the historical
relationship between the languages: Python contains many C-originated security idioms that
were adopted as the language incorporated system calls, while C does not contain the
Python-specific idioms.
This directional asymmetry suggests that pre-training corpus properties --- the frequency
and distribution of cross-language overlapping tokens --- play a meaningful role in
determining transfer direction.

\section{The Synthetic Data Ceiling}

The residual gap of 0.182 F1 points between synthetic and real-world performance, even
after augmentation, reflects a fundamental limitation of synthetic training data.
Hand-crafted training examples are, by construction, unambiguous: each vulnerable example
contains exactly one exploitable pattern presented in its most recognisable form.
Real-world functions are more complex in several ways.

First, real functions are longer and contain multiple operations, some of which superficially
resemble vulnerability patterns without being exploitable.
A function that performs bounds checking before an array access contains the same
array-indexing token sequence as an out-of-bounds vulnerability; distinguishing them
requires understanding the guard condition, not just the access itself.
Second, real vulnerabilities are often latent: the dangerous operation arises from the
interaction between multiple lines across a function body.
Third, real code contains domain-specific idioms and coding styles that differ from the
hand-crafted training examples in ways the model has never been exposed to.

The augmentation with 1{,}000 DiverseVul samples partially addresses the third issue
by exposing the model to real code styles, but does not resolve the first two issues ---
they require architectural changes rather than more training data.
The asymptote of the data augmentation approach is the ceiling imposed by the model's
inability to reason about inter-statement program properties.

A related limitation is class coverage.
The real-world test set contains CWE-787, CWE-400, CWE-703, and CWE-284, none of which
appear in the synthetic training set.
Recall on unseen CWE classes is zero.
Expanding the synthetic dataset to cover additional classes, or incorporating a larger
real-world training corpus --- DiverseVul provides approximately 200{,}000 labelled
C functions; only 1{,}000 are used here --- would directly address this gap.

\section{Calibration and Practical Deployment}

The near-unity temperature ($T = 0.981$) after calibration indicates that the
label-smoothed model is well-calibrated on synthetic data without requiring post-hoc
adjustment.
Model outputs can be treated as approximate probabilities and used to prioritise a
vulnerability report queue.

However, calibration was assessed on synthetic data.
The real-world calibration plot shows a broader, less bimodal distribution than the
synthetic one, reflecting the model's uncertainty about real-world functions that partially
match training patterns.
A deployment-quality calibration assessment would require a substantially larger held-out
real-world set spanning the same CWE distribution as the deployment target.

Threshold selection is a further practical consideration.
F1 varies relatively little between thresholds 0.1 and 0.7 (0.928 to 0.927), but the
precision/recall trade-off within that range is substantial.
An organisation using this system for automated code review has a different operating
requirement than one using it for manual audit support.
The threshold sweep in Experiment~4 provides the data to make this choice explicitly,
which is more informative than reporting a single threshold result.

\section{Limitations Summary}

The principal limitations of the system are:

\begin{enumerate}
  \item \textbf{Architectural.}
        The model cannot reason about control-flow reachability (dead code) or data-flow
        provenance (whether a variable is user-controlled).
        Failures on CWE-78 C minimal-fix and all dead-code test cases follow directly from this.

  \item \textbf{CWE coverage.}
        The synthetic training set covers 17 CWE types.
        The DiverseVul test set contains approximately 50 distinct CWE types.
        Recall on unseen CWE classes is zero.

  \item \textbf{Language coverage.}
        The model is trained on C and Python only.
        Applying it to Java, JavaScript, or Go would require retraining or fine-tuning
        on data from those languages.

  \item \textbf{Synthetic distribution shift.}
        Even with augmentation, the training distribution is skewed towards clean,
        minimal, function-isolated vulnerability patterns.
        Real-world functions are longer, noisier, and more context-dependent.

  \item \textbf{Real-world evaluation scale.}
        The real-world test set contains 200 samples.
        This is sufficient to measure overall accuracy but insufficient to estimate
        per-CWE recall reliably for low-frequency classes.

  \item \textbf{Single-language real-world data.}
        The DiverseVul augmentation is exclusively C/C++.
        A Python vulnerability dataset analogous to DiverseVul at comparable scale does
        not currently exist.
\end{enumerate}

